% !TEX program = xelatex
% Заготовка ВКР по требованиям ИЕНиМ (2024) и ГОСТ 7.32–2017.
% Запуск: latexmk -xelatex -interaction=nonstopmode thesis.tex
\documentclass[a4paper,14pt]{extarticle}

\usepackage{fontspec}
\setmainfont{Times New Roman}
\usepackage{polyglossia}
\setdefaultlanguage{russian}
\newfontfamily\cyrillicfont{Times New Roman}[Script=Cyrillic]
\newfontfamily\cyrillicfontsf{Times New Roman}[Script=Cyrillic]
\newfontfamily\cyrillicfonttt{Liberation Mono}

\usepackage[left=30mm,right=15mm,top=20mm,bottom=20mm]{geometry}
\usepackage{setspace}
\usepackage{array}
\usepackage{booktabs}
\usepackage{graphicx}
\usepackage{amsmath}
\usepackage{unicode-math}
\setmathfont{TeX Gyre Termes Math}
\usepackage{caption}
\usepackage{titlesec}
\usepackage[hidelinks]{hyperref}

\onehalfspacing
\setlength{\parindent}{1.25cm}
\setlength{\parskip}{0pt}
\setlength{\emergencystretch}{3em}
\tolerance=2000
\setcounter{secnumdepth}{3}
\setcounter{tocdepth}{3}

% Заголовки: структурные элементы — по центру, разделы — с абзаца.
% Переносы в заголовках запрещены (ИЕНиМ, с. 13; УрФУ, с. 40).
\titleformat{\section}[block]{\bfseries\normalsize\hyphenpenalty=10000\exhyphenpenalty=10000}{\thesection}{0.5em}{}
\titlespacing*{\section}{1.25cm}{1.5\baselineskip}{1\baselineskip}
\titleformat{\subsection}[block]{\bfseries\normalsize\hyphenpenalty=10000\exhyphenpenalty=10000}{\thesubsection}{0.5em}{}
\titlespacing*{\subsection}{1.25cm}{1\baselineskip}{0.5\baselineskip}
\titleformat{\subsubsection}[block]{\bfseries\normalsize\hyphenpenalty=10000\exhyphenpenalty=10000}{\thesubsubsection}{0.5em}{}
\titlespacing*{\subsubsection}{1.25cm}{0.75\baselineskip}{0.25\baselineskip}

% Структурный заголовок: прописными, полужирным, по центру строки.
% Отступ 1,25 см здесь не применяется — только к разделам и подразделам.
\makeatletter
\newcommand{\structheading}[1]{%
  \phantomsection
  \begingroup
    \centering
    \parindent=0pt
    \hyphenpenalty=10000\exhyphenpenalty=10000
    \textbf{#1}\par
  \endgroup
  \vspace{0.5\baselineskip}}
\newcommand{\structtitle}[1]{%
  \clearpage
  \structheading{#1}%
  \addcontentsline{toc}{section}{#1}}
\newcommand{\fronttitle}[1]{%
  \clearpage
  \structheading{#1}}
\newcommand{\vkrtableofcontents}{\@starttoc{toc}}
% Отточие после каждого элемента содержания (ИЕНиМ, с. 10; УрФУ, с. 24).
\renewcommand{\l@section}{\@dottedtocline{1}{0pt}{1.5em}}
\makeatother

% Подписи: 12 pt, разделитель — тире (ИЕНиМ, с. 14–16; УрФУ, с. 49).
% Подпись рисунка: одна строка — по центру, несколько — по ширине.
% Название таблицы: слева, без абзацного отступа.
\DeclareCaptionFont{vkrcaption}{\small}
\captionsetup{font=vkrcaption,labelsep=endash,justification=justified,singlelinecheck=true}
\captionsetup[table]{justification=raggedright,singlelinecheck=false,indention=0pt}

\pagestyle{plain}
\makeatletter
\renewcommand{\@oddfoot}{\hfil\thepage\hfil}
\renewcommand{\@evenfoot}{\hfil\thepage\hfil}
\makeatother

% Нумерация рисунков, таблиц и формул — по разделу: Рисунок 2.1, Таблица 3.1, (4.1).
% ИЕНиМ, с. 14–16. Для приложения А:
% \renewcommand{\thefigure}{А.\arabic{figure}}\setcounter{figure}{0}
\counterwithin{figure}{section}
\counterwithin{table}{section}
\counterwithin{equation}{section}

% ИЕНиМ, с. 16: выше и ниже выключной формулы — не менее одной свободной строки.
\AtBeginDocument{%
  \setlength{\abovedisplayskip}{\baselineskip}%
  \setlength{\belowdisplayskip}{\baselineskip}%
  \setlength{\abovedisplayshortskip}{\baselineskip}%
  \setlength{\belowdisplayshortskip}{\baselineskip}}

% Список источников: арабские цифры без точки, с абзацного отступа (ИЕНиМ, с. 18).
\newcounter{source}
\newenvironment{sourcelist}{%
  \begin{list}{\arabic{source}}{%
    \usecounter{source}%
    \setlength{\leftmargin}{\dimexpr\parindent+1.2em\relax}%
    \setlength{\labelwidth}{0.7em}%
    \setlength{\labelsep}{0.5em}%
    \setlength{\itemindent}{0pt}%
    \setlength{\itemsep}{0pt}%
    \setlength{\parsep}{0pt}%
    \setlength{\topsep}{0pt}%
    \setlength{\partopsep}{0pt}}%
}{\end{list}}

% Строки реферата: подчёркивание не выходит за правое поле.
\newlength{\refw}
\newcommand{\refblankline}[1]{%
  \par
  \noindent\hspace*{\parindent}%
  \settowidth{\refw}{#1}%
  #1\underline{\hspace{\dimexpr\linewidth-\parindent-\refw-0.5em\relax}}\par}
\newcommand{\refkeywords}[1]{%
  \par
  \noindent
  \settowidth{\refw}{#1}%
  #1\ \underline{\hspace{\dimexpr\linewidth-\refw-0.5em\relax}}\par}

% —————————————————————————
% Данные работы (замените значения перед началом набора текста)
% Эти макросы выводятся и на титульном листе, и в задании, и в реферате.
% —————————————————————————
\newcommand{\VKRTitle}{ХРОМАТОГРАФИЧЕСКОЕ ИССЛЕДОВАНИЕ СИСТЕМЫ ПОЛИЭТИЛЕНОКСИД–ХИТОЗАН}
\newcommand{\VKRStudent}{Смирнов Юрий Викторович}
\newcommand{\VKRStudentGenitive}{Смирнова Юрия Викторовича}
% Для титула Ф. И. О. разбивается на две строки, как в образце ИЕНиМ (с. 24).
\newcommand{\VKRStudentGenitiveTitle}{Смирнова\\Юрия Викторовича}
\newcommand{\VKRGruppa}{\underline{\hspace{2.5cm}}}
\newcommand{\VKRDirection}{04.03.01 «Химия»}
% Образовательная программа указывается при наличии нескольких программ
% направления подготовки (ИЕНиМ, с. 24, сноска 11).
\newcommand{\VKRProgram}{«Химия»}
\newcommand{\VKRInstitute}{Институт естественных наук и математики}
\newcommand{\VKRInstituteLoc}{Институте естественных наук и математики}
\newcommand{\VKRDepartment}{Кафедра органической химии и высокомолекулярных соединений}
\newcommand{\VKRDepartmentLoc}{кафедре органической химии и высокомолекулярных соединений}
\newcommand{\VKRHeadPosition}{Зав.~кафедрой}
\newcommand{\VKRHeadDegree}{д.~х.~н., проф.}
\newcommand{\VKRHeadName}{В.~Я.~Сосновских}
\newcommand{\VKRNormControlDegree}{к.~х.~н., доц.}
\newcommand{\VKRNormControlName}{Л.~В.~Адамова}
\newcommand{\VKRAdvisorDegree}{к.~х.~н., доц.}
\newcommand{\VKRAdvisor}{И.~С.~Иванов}
\newcommand{\VKRWorkType}{Выпускная квалификационная работа бакалавра}
\newcommand{\VKRYear}{2026}
\newcommand{\VKRPlace}{Екатеринбург}

\hypersetup{
  pdftitle={\VKRTitle},
  pdfauthor={\VKRStudent},
  pdfsubject={Выпускная квалификационная работа}
}

\begin{document}

% polyglossia выбирает язык при \begin{document} и перезаписывает названия
% элементов, поэтому задаём их здесь, после выбора языка.
\renewcommand{\contentsname}{СОДЕРЖАНИЕ}
\renewcommand{\refname}{СПИСОК ИСПОЛЬЗОВАННЫХ ИСТОЧНИКОВ И ЛИТЕРАТУРЫ}
\renewcommand{\figurename}{Рисунок}
\renewcommand{\tablename}{Таблица}

% =============================== ТИТУЛЬНЫЙ ЛИСТ =============================
% Векторный титульный лист по образцу ИЕНиМ (2024), с. 24–26.
% Поля и шрифт наследуются от основного документа; титул набирается
% одинарным интервалом, как в образце (основной текст — 1,5).
% Грифы СОГЛАСОВАНО/УТВЕРЖДАЮ не печатаются: в образце ИЕНиМ их нет.
\thispagestyle{empty}
\begingroup
\singlespacing
\setlength{\parindent}{0pt}
\hyphenpenalty=10000\exhyphenpenalty=10000

{\centering
МИНИСТЕРСТВО НАУКИ И ВЫСШЕГО ОБРАЗОВАНИЯ РОССИЙСКОЙ ФЕДЕРАЦИИ\\
федеральное государственное автономное образовательное учреждение\\
высшего образования\\
УРАЛЬСКИЙ ФЕДЕРАЛЬНЫЙ УНИВЕРСИТЕТ\\
имени первого Президента России Б.~Н.~Ельцина\par}

\vspace{1.5\baselineskip}

{\centering \MakeUppercase{\VKRInstitute}\par}

\vspace{3.5\baselineskip}

{\centering \VKRDepartment\par}

\vspace{2.5\baselineskip}

{\centering \textbf{\VKRTitle}\par}

\vspace{1.5\baselineskip}

{\centering
Направление подготовки \VKRDirection\\
Образовательная программа \VKRProgram\par}

\vspace{3\baselineskip}

\noindent
\begin{minipage}[t]{0.48\textwidth}
\VKRHeadPosition:\\[0.2em]
\VKRHeadDegree\ \VKRHeadName\\[2.2em]
\rule{5.5cm}{0.4pt}
\end{minipage}\hfill
\begin{minipage}[t]{0.48\textwidth}
\centering
\VKRWorkType\\[1em]
\textbf{\VKRStudentGenitiveTitle}\\[1.2em]
\rule{5.5cm}{0.4pt}
\end{minipage}

\vspace{2\baselineskip}

\noindent
\begin{minipage}[t]{0.48\textwidth}
Нормоконтролер:\\[0.2em]
\VKRNormControlDegree\ \VKRNormControlName\\[2.2em]
\rule{4.5cm}{0.4pt}
\end{minipage}\hfill
\begin{minipage}[t]{0.48\textwidth}
Научный руководитель:\\[0.2em]
\VKRAdvisorDegree\ \VKRAdvisor\\[2.2em]
\rule{5.5cm}{0.4pt}
\end{minipage}

\vfill

{\centering
\VKRPlace\\
\VKRYear\par}
\endgroup
\clearpage

% Титульный лист учитывается в общей нумерации, но цифра на нём не печатается.
\setcounter{page}{2}

% ================================ ЗАДАНИЕ ===================================
% Бланк задания для печати и подписания. Подписанный скан можно вставить
% вместо этой страницы, сохранив номер страницы 2.

\begin{flushright}
УТВЕРЖДАЮ\\[0.2em]
Руководитель образовательной программы\\[0.5em]
\underline{\hspace{4cm}}\quad(\underline{\hspace{2.6cm}})\\[0.3em]
«\underline{\hspace{0.8cm}}» \underline{\hspace{3cm}} 20\underline{\hspace{0.5cm}} г.
\end{flushright}

\vspace{0.5em}
\noindent Код, наименование направления: \VKRDirection\\
\noindent Наименование программы: \underline{\hspace{0.55\linewidth}}\\
\noindent Группа: \VKRGruppa

\begin{center}
\textbf{ЗАДАНИЕ}\\
на выполнение выпускной квалификационной работы
\end{center}

\noindent студента \VKRStudentGenitive\\
(Фамилия, имя, отчество)

\noindent Квалификация: бакалавр\\
(бакалавр, специалист, магистр)

{\hyphenpenalty=10000\exhyphenpenalty=10000
\noindent Провести научное исследование по теме: \VKRTitle\par}

\begin{center}
Срок представления работы научному руководителю:\\
«\underline{\hspace{0.8cm}}» \underline{\hspace{1.6cm}} 20\underline{\hspace{0.55cm}} г.
\end{center}

\noindent Научный руководитель \hfill \underline{\hspace{4.5cm}}\quad(Подпись)
\\(\VKRAdvisor)

\vspace{0.5em}
\noindent Задание принял к исполнению \hfill \underline{\hspace{3.5cm}}\\
(Подпись)

% =============================== РЕФЕРАТ ====================================
\fronttitle{РЕФЕРАТ}
Отчет \underline{\hspace{1.2cm}} с., \underline{\hspace{0.8cm}} рис., \underline{\hspace{0.8cm}} табл., \underline{\hspace{1.2cm}} источн., \underline{\hspace{0.8cm}} прил.

\refkeywords{Ключевые слова:}
\refblankline{Объект исследования~— }
\refblankline{Цель работы~— }
\refblankline{Методы исследования~— }
\refblankline{Результаты и их новизна~— }
\refblankline{Область применения~— }
\refblankline{Значимость работы~— }

\fronttitle{МЕСТО ВЫПОЛНЕНИЯ РАБОТЫ}
Работа выполнена в \VKRInstituteLoc, на \VKRDepartmentLoc{} Уральского федерального университета, г.~\VKRPlace.

% =============================== СОДЕРЖАНИЕ =================================
\fronttitle{\contentsname}
\vkrtableofcontents

% ========================= ОБОЗНАЧЕНИЯ И СОКРАЩЕНИЯ =========================
\structtitle{ОБОЗНАЧЕНИЯ И СОКРАЩЕНИЯ}
В настоящей работе применяют следующие обозначения и сокращения.

\noindent\begin{tabular}{@{}>{\raggedright\arraybackslash}p{0.22\textwidth}>{\raggedright\arraybackslash}p{0.73\textwidth}@{}}
ВКР & — выпускная квалификационная работа\\
Добавьте & — расшифровка сокращения
\end{tabular}

% ================================ ВВЕДЕНИЕ ==================================
\structtitle{ВВЕДЕНИЕ}
Во введении необходимо обосновать актуальность темы, сформулировать объект и предмет исследования, цель и задачи работы, методы исследования, научную новизну, практическую значимость и положения, выносимые на защиту.

% =============================== ОСНОВНАЯ ЧАСТЬ =============================
\clearpage\section{Обзор литературы}
Текст подраздела. Ссылки на источники оформляются в квадратных скобках, например [1].

\clearpage\section{Постановка задачи}
Сформулируйте исследовательскую задачу и критерии оценки результатов.

\clearpage\section{Методика эксперимента}
\subsection{Материалы и оборудование}
Опишите объект исследования, материалы, приборы и условия эксперимента.
\subsection{Методика проведения исследования}
Опишите последовательность теоретических и экспериментальных процедур.

\clearpage\section{Результаты и их обсуждение}
\subsection{Результаты}
Представьте результаты с необходимыми рисунками, таблицами и расчетами.

% Пример оформления рисунка — замените изображение и подпись своим материалом.
\begin{figure}[htbp]
  \centering
  \fbox{\rule{0pt}{3cm}\rule{8cm}{0pt}}
  \caption{Пример оформления рисунка}
  \label{fig:example}
\end{figure}

На рисунке~\ref{fig:example} приведен пример оформления иллюстрации.

% Пример оформления таблицы — замените данные своим материалом.
В таблице~\ref{tab:example} приведен пример оформления табличных данных.

\begin{table}[htbp]
  \caption{Пример оформления таблицы}
  \label{tab:example}
  \begin{tabular}{@{}lcc@{}}
    \toprule
    Показатель & Значение 1 & Значение 2 \\
    \midrule
    Пример & 1,0 & 2,0 \\
    Пример & 3,0 & 4,0 \\
    \bottomrule
  \end{tabular}
\end{table}

% Пример оформления формулы — замените своим материалом.
Уравнение состояния идеального газа имеет вид
\begin{equation}
  pV = nRT,
  \label{eq:state}
\end{equation}

где $p$ — давление; $V$ — объем; $n$ — количество вещества; $R$ — универсальная газовая постоянная; $T$ — температура.

Формула~\eqref{eq:state} приведена в качестве примера оформления.

\subsection{Обсуждение результатов}
Сопоставьте результаты с опубликованными данными и оцените достоверность.
Дополнительные материалы, необходимые для понимания результатов, приведены в приложении~А.

\structtitle{ЗАКЛЮЧЕНИЕ (ВЫВОДЫ)}
Сформулируйте краткие выводы по задачам, оцените полноту решения, практическую значимость и возможные направления дальнейшей работы.

\structtitle{\refname}
\begin{sloppypar}
\begin{sourcelist}
  \item Автор. Название источника. – Место издания : Издательство, год. – 100 с.
  \item Автор статьи. Название статьи // Журнал. – 2024. – Т. 35, № 1. – С. 1–10.
\end{sourcelist}
\end{sloppypar}

% =============================== ПРИЛОЖЕНИЕ =================================
\clearpage
\phantomsection
\addcontentsline{toc}{section}{ПРИЛОЖЕНИЕ А (обязательное). ДОПОЛНИТЕЛЬНЫЕ МАТЕРИАЛЫ}
\begingroup
  \centering
  \parindent=0pt
  \hyphenpenalty=10000\exhyphenpenalty=10000
  \textbf{ПРИЛОЖЕНИЕ А}\par
  (обязательное)\par
  \textbf{ДОПОЛНИТЕЛЬНЫЕ МАТЕРИАЛЫ}\par
\endgroup
\vspace{0.5\baselineskip}

Материалы приложения, необходимые для понимания результатов исследования. Каждое следующее приложение начинайте с новой страницы.

\end{document}
